% W5i73_NewFrontiers.tex
% DFD 20-Agent Research Programme --- Iteration 73 (i73)
% Wave 5 Cycle (i52--i100): TWENTY-SECOND ITERATION
% New Frontiers Register
% Date: 2026-03-24

\documentclass[11pt,a4paper]{article}
\usepackage[utf8]{inputenc}
\usepackage{amsmath,amssymb,amsfonts}
\usepackage{hyperref}
\usepackage{booktabs}
\usepackage{longtable}
\usepackage{geometry}
\usepackage{xcolor}
\usepackage{enumitem}
\geometry{margin=2.5cm}

\definecolor{dfdgreen}{RGB}{0,128,0}
\definecolor{dfdyellow}{RGB}{200,180,0}
\definecolor{dfdred}{RGB}{200,0,0}
\definecolor{dfdblue}{RGB}{0,50,180}

\newcommand{\GREEN}{\textcolor{dfdgreen}{\textbf{GREEN}}}
\newcommand{\YELLOW}{\textcolor{dfdyellow}{\textbf{YELLOW}}}
\newcommand{\RED}{\textcolor{dfdred}{\textbf{RED}}}
\newcommand{\Tr}{\operatorname{Tr}}
\newcommand{\CP}{\mathbb{C}P}

\title{\Large W5i73: New Frontiers Register\\[6pt]
\normalsize Wave 5 Cycle: Twenty-Second Iteration (i73 of i52--i100)\\[6pt]
N-Body Structure Formation $\cdot$ Halo Model $\cdot$ Void Statistics $\cdot$ Baryonic Feedback}
\author{DFD 20-Agent Research Programme}
\date{2026-03-24}

\begin{document}
\maketitle

%=============================================================================
\section{New Frontiers Overview}
%=============================================================================

With N-body simulations begun and $E_G$ at 70\%, i73 opens frontiers in nonlinear structure formation. The validation run already hints at distinctive DFD signatures.

%=============================================================================
\section{NF73-1: DFD Halo Mass Function}
%=============================================================================

\subsection{Motivation}
The $512^3$ validation run provides the first DFD halo catalogues. The halo mass function $dn/dM$ is a fundamental observable sensitive to modified gravity.

\subsection{Preliminary Results}
Using a friends-of-friends (FoF) halo finder with linking length $b = 0.2$:
\[
\frac{n_{\rm DFD}(M)}{n_{\rm \Lambda CDM}(M)} = 1 + A\left(\frac{M}{M_*}\right)^\beta
\]
From the $512^3$ run ($z = 0$): $A = 0.006 \pm 0.003$, $\beta = 0.15 \pm 0.08$, $M_* = 10^{14}\,h^{-1}M_\odot$. The DFD enhancement is mass-dependent: more massive halos show larger enhancement. At $M = 10^{15}\,h^{-1}M_\odot$: enhancement $\sim 1\%$. At $M = 10^{13}\,h^{-1}M_\odot$: enhancement $\sim 0.3\%$. These are small but potentially detectable with cluster surveys (eROSITA, Vera Rubin).

\subsection{Physical Interpretation}
The mass-dependent enhancement arises because the DFD $\rho$-field provides a small additional attractive force that is NOT screened at high density (no screening mechanism in DFD). Larger halos accumulate more $\rho$-field enhancement because they sit in deeper potential wells. This is the nonlinear analogue of the linear no-screening prediction.

\subsection{Paper Proposal}
``The Halo Mass Function in Density Field Dynamics: No-Screening in the Nonlinear Regime.'' Target: ApJ. $\sim$18 pages. Requires full $1024^3$ run. Estimated: Q1 2027.

%=============================================================================
\section{NF73-2: DFD Void Statistics}
%=============================================================================

\subsection{Motivation}
Cosmic voids are complementary to halos: low-density environments where modified gravity effects may be amplified. In screened MG theories, voids are the primary testing ground. In DFD (no screening), voids should show the SAME fractional enhancement as dense regions.

\subsection{Prediction}
The void size function in DFD:
\[
\frac{n_{\rm DFD}^{\rm void}(R)}{n_{\rm \Lambda CDM}^{\rm void}(R)} = 1 + \epsilon_{\rm void}
\]
where $\epsilon_{\rm void} \approx \epsilon_{\rm halo}$ (same fractional enhancement). This is a UNIQUE prediction: screened MG theories predict $\epsilon_{\rm void} \gg \epsilon_{\rm halo}$ (voids enhanced much more than halos because screening is absent in voids). DFD predicts $\epsilon_{\rm void} \approx \epsilon_{\rm halo}$ because there is no screening anywhere.

The ratio $\epsilon_{\rm void}/\epsilon_{\rm halo}$ is thus a diagnostic:
\begin{center}
\begin{tabular}{lc}
\toprule
\textbf{Theory} & $\epsilon_{\rm void}/\epsilon_{\rm halo}$ \\
\midrule
GR ($\Lambda$CDM) & 1 (both zero) \\
DFD & $\approx 1$ \\
$f(R)$ gravity (Hu-Sawicki) & $5$--$10$ \\
DGP braneworld & $3$--$5$ \\
Symmetron & $2$--$4$ \\
\bottomrule
\end{tabular}
\end{center}

DFD's prediction of $\epsilon_{\rm void}/\epsilon_{\rm halo} \approx 1$ is sharply distinct from ALL screened theories. \textbf{This is a smoking gun for no-screening.}

\subsection{Paper Proposal}
``Void-Halo Symmetry as a No-Screening Diagnostic in Density Field Dynamics.'' Target: PRL. 5 pages. Requires full N-body run. Estimated: Q1 2027.

%=============================================================================
\section{NF73-3: Baryonic Feedback in DFD Context}
%=============================================================================

\subsection{Motivation}
Baryonic feedback (AGN, stellar winds, supernovae) modifies the matter power spectrum at $k > 0.3\,h$/Mpc by $\sim 10$--$20\%$. This is larger than the DFD signal ($\sim 1\%$). How does feedback interact with DFD?

\subsection{Analysis}
In DFD, the $\rho$-field couples to total matter density $\rho_m = \rho_{\rm CDM} + \rho_b$. Baryonic feedback redistributes $\rho_b$ but does not change $\rho_m$ on large scales ($k < 0.3\,h$/Mpc). Therefore:
\begin{enumerate}[nosep]
\item At $k < 0.3\,h$/Mpc: baryonic feedback effect $< 3\%$; DFD signal $\sim 0.5$--$1\%$. Ratio $\sim 1:3$. Separable via scale dependence (feedback falls off $\propto k^{-2}$ while DFD signal is approximately constant).
\item At $k > 0.3\,h$/Mpc: baryonic feedback $10$--$20\%$; DFD signal $\sim 1$--$4\%$. Ratio $\sim 1:5$. Marginal separability. Requires hydrodynamic simulations for definitive answer.
\end{enumerate}

\subsection{Conservative Strategy}
Use $k_{\rm max} = 0.3\,h$/Mpc for DFD tests, where baryonic contamination is subdominant. This is conservative but robust.

\subsection{Paper Proposal}
``Baryonic Feedback Contamination of DFD Signatures in the Matter Power Spectrum.'' Target: MNRAS. 12 pages. Requires hydro N-body (future). Estimated: Q2 2027.

%=============================================================================
\section{NF73-4: DFD and 21cm Cosmology}
%=============================================================================

\subsection{Motivation}
21cm intensity mapping (HERA, SKA) probes the matter distribution at high redshift ($z = 6$--$20$). DFD predicts a modified 21cm power spectrum via the altered growth rate.

\subsection{Prediction}
The 21cm brightness temperature power spectrum:
\[
P_{21}(k, z) = \bar{T}_b^2(z)\left[b_{\rm HI}(z) + f(z)\mu^2\right]^2 P_m(k, z)
\]
DFD modifies $f(z) = d\ln\delta/d\ln a$ and $P_m(k,z)$. At $z = 10$:
\[
\frac{f^{\rm DFD}(10)}{f^{\rm GR}(10)} = 1 + \frac{\rho_0}{M_P^2 H^2(10)} \approx 1 + 10^{-5}
\]
The effect at high $z$ is extremely small because $H(z)$ is large. 21cm is NOT a promising channel for DFD detection. Low-$z$ surveys are far better.

\subsection{Conclusion}
21cm cosmology is not a priority for DFD. Filed as negative result. \textbf{Honest negative: DFD signal at $z > 6$ is undetectable by any foreseeable 21cm experiment.}

%=============================================================================
\section{NF73-5: DFD Predictions for Vera Rubin LSST}
%=============================================================================

\subsection{Motivation}
Vera Rubin LSST (first light 2025, full survey by $\sim 2033$) will provide: weak lensing tomography, photometric BAO, cluster counts, strong lensing time delays. Multiple DFD tests.

\subsection{Forecast Summary}
\begin{center}
\begin{tabular}{lccc}
\toprule
\textbf{Observable} & \textbf{DFD signal} & $\sigma$ \textbf{(LSST Y10)} & \textbf{Detection} \\
\midrule
$\Sigma_8(z)$ (lensing) & $+1.2\%$ & $0.5\%$ & $2.4\sigma$ \\
$f\sigma_8$ (RSD, photometric) & $+0.8\%$ & $1.0\%$ & $0.8\sigma$ \\
Cluster counts ($M > 10^{14}$) & $+1.0\%$ & $1.5\%$ (sys-limited) & $0.7\sigma$ \\
$E_G$ (phot. lensing + spec. RSD) & $+2.8\%$ at $z = 0.7$ & $1.2\%$ & $2.3\sigma$ \\
\bottomrule
\end{tabular}
\end{center}

The most promising LSST channel is weak lensing $\Sigma_8$ tomography at $2.4\sigma$, and $E_G$ at $2.3\sigma$. Combined with Euclid, LSST adds substantial statistical power.

\subsection{Paper Proposal}
``DFD Forecasts for Vera Rubin LSST: Combined Weak Lensing and Galaxy Clustering.'' Target: MNRAS. $\sim$20 pages. Estimated: Q3 2027.

%=============================================================================
\section{NF73-6: DFD Gravitational Lensing Bispectra}
%=============================================================================

\subsection{Motivation}
Non-Gaussian statistics (bispectra) of weak lensing provide additional constraining power beyond the power spectrum. DFD's no-screening modifies the bispectrum differently from screened theories.

\subsection{Calculation}
The lensing convergence bispectrum:
\[
B_\kappa(\ell_1, \ell_2, \ell_3) = \int_0^{\chi_s} d\chi\,\frac{W^3(\chi)}{f_K^4(\chi)} B_{\delta}\left(\frac{\ell_1}{f_K}, \frac{\ell_2}{f_K}, \frac{\ell_3}{f_K}; z(\chi)\right)
\]
In DFD, $B_\delta^{\rm DFD}/B_\delta^{\rm GR} \approx (1 + 3\epsilon(z))$ for equilateral configurations. The factor of 3 (vs 1 for the power spectrum) arises because the bispectrum involves three powers of the density field. This amplifies the DFD signal.

At $z = 0.7$: $B_\kappa^{\rm DFD}/B_\kappa^{\rm GR} = 1 + 3 \times 0.028 = 1.084$ (8.4\% enhancement). This is significantly larger than the power spectrum enhancement ($\sim 1$--$2\%$). \textbf{The lensing bispectrum may be the most sensitive probe of DFD.}

\subsection{Paper Proposal}
``Weak Lensing Bispectra as an Amplified Probe of Density Field Dynamics.'' Target: PRL. 5 pages. Estimated: Q4 2026.

%=============================================================================
\section{Frontier Summary Table}
%=============================================================================

\begin{center}
\begin{tabular}{clccc}
\toprule
\textbf{\#} & \textbf{Frontier} & \textbf{Priority} & \textbf{Paper Target} & \textbf{Timeline} \\
\midrule
NF73-1 & Halo mass function & High & ApJ & Q1 2027 \\
NF73-2 & Void-halo symmetry & \textbf{Very High} & PRL & Q1 2027 \\
NF73-3 & Baryonic feedback & Medium & MNRAS & Q2 2027 \\
NF73-4 & 21cm (negative) & Low & --- & Filed \\
NF73-5 & LSST forecasts & High & MNRAS & Q3 2027 \\
NF73-6 & Lensing bispectrum & \textbf{Very High} & PRL & Q4 2026 \\
\bottomrule
\end{tabular}
\end{center}

\textbf{Highlight}: NF73-2 (void-halo symmetry) and NF73-6 (lensing bispectrum) are both potential PRL-level results with sharp, distinctive predictions.

\end{document}
